\documentclass[12pt]{article}
\usepackage{graphicx}
\newcommand{\maybeincludegraphics}[2][]{%
  \IfFileExists{#2}{%
    \includegraphics[#1]{#2}%
  }{%
    \fbox{\parbox{0.95\linewidth}{\centering Missing figure: \texttt{\detokenize{#2}}}}%
  }%
}
\usepackage{natbib}
\bibliographystyle{agsm}
\usepackage{tabularx}
\usepackage[table]{xcolor}
\usepackage{colortbl}
\usepackage{mathtools}
\usepackage{booktabs}
\usepackage{multirow}
\usepackage{array}
\usepackage{ragged2e}
\usepackage{caption}
\captionsetup[table]{justification=raggedright,singlelinecheck=false}
\usepackage{amsmath}
\usepackage{amssymb}
\usepackage[hyperfootnotes=false]{hyperref}
\hypersetup{linkcolor=black,urlcolor=blue}
\usepackage{lmodern}
\usepackage[nopatch=footnote]{microtype}
\usepackage{pdflscape}
\usepackage{newtxtext,newtxmath}
\usepackage[a4paper, margin=2.5cm]{geometry}
\usepackage{setspace}
\singlespacing
\usepackage[T1]{fontenc}
\usepackage[utf8]{inputenc}
\usepackage{titlesec}

\titleformat{\section}{\bfseries\large}{\thesection}{1em}{}
\titleformat{\subsection}{\bfseries\itshape\normalsize}{\thesubsection}{1em}{}
\titleformat{\subsubsection}{\itshape\normalsize}{\thesubsubsection}{1em}{}

\begin{document}

\title{\textbf{Liquidity Premia, Arbitrage Bands, and the Collapse of Put--Call Parity Violations in Cryptocurrency Option Markets}}

\author{Balázs Králik\thanks{Corresponding Author. Assistant Professor at Corvinus University of Budapest, Institute of Finance, Fővám square 8. Budapest, Hungary, 1093. E-mail: balazs.kralik@uni-corvinus.hu.}, Nóra Felföldi-Szűcs\thanks{Associate Professor at Corvinus University of Budapest, Institute of Finance.}, Tamás Rutai\thanks{Student at Corvinus University of Budapest.}, Kata Váradi\thanks{Associate Professor at Corvinus University of Budapest, Institute of Finance.}}
\date{}

\maketitle

\begin{abstract}
This paper studies put--call parity deviations in cryptocurrency option markets as a form of liquidity premium rather than as a simple violation of a frictionless pricing identity. We build on the Absolute Market Unfairness (AMU) framework, which measures the risk-free arbitrage profits implied by quoted bid--ask prices, and reinterpret AMU as the positive part of a no-arbitrage wedge generated by transaction costs, funding constraints, collateral segmentation, legal uncertainty, and limited arbitrage capital. The starting point is the empirical evidence from OKX and Deribit BTCUSD and ETHUSD option markets that AMU is concentrated in less liquid segments of the option surface: ETH relative to BTC, away-from-the-money options relative to at-the-money options, short-dated contracts relative to longer-dated contracts, and periods of market stress relative to calmer periods. The broader empirical motivation is the apparent collapse of put--call-parity-breaking opportunities around the institutionalization of crypto markets in 2024, a period marked by the approval of U.S. spot Bitcoin exchange-traded products, subsequent spot Ether ETF approvals, and the staged implementation of MiCA in the European Union. We propose a threshold-arbitrage model in which parity deviations are not eliminated unless they exceed an endogenous arbitrage band. Regulatory clarity, better custody, exchange competition, and deeper arbitrage capital narrow this band, producing a measurable transition from visible but non-bankable arbitrage to institutionally enforceable arbitrage. The paper positions crypto option parity violations alongside classical historical examples of liquidity premia and limits to arbitrage, including covered-interest-parity deviations, ETF premium/discount arbitrage, on-the-run/off-the-run Treasury spreads, and CDS--bond basis trades. We outline an empirical strategy using high-frequency Deribit and OKX option and futures quotes from 2020--2026, exchange-level order-book liquidity, funding rates, stablecoin basis measures, ETF-flow variables, and regulatory-event indicators. The resulting framework converts AMU from a cross-sectional market-quality statistic into a dynamic measure of crypto-market integration and liquidity-premium compression.

\medskip
\textbf{Keywords:} cryptocurrency options, put-call parity, liquidity premium, limits to arbitrage, market efficiency, ETF approval, MiCA

\medskip
\textbf{JEL classification:} G12, G14, G15, G23
\end{abstract}

\section{Introduction}

Put--call parity (PCP) is one of the most elementary no-arbitrage restrictions in option pricing. In a frictionless market, a European call and put with the same strike and maturity must differ in value by the forward value of the underlying. Yet in actual markets parity is not a point restriction. It is a band. Trading costs, bid--ask spreads, funding rates, margin requirements, collateral constraints, short-sale costs, settlement risk, and legal uncertainty determine whether an apparent arbitrage can actually be implemented.

This paper studies PCP deviations in cryptocurrency option markets through this liquidity-premium lens. We build on the Absolute Market Unfairness (AMU) measure developed in our companion paper, where AMU is defined as the risk-free arbitrage profit implied by quoted option and futures bid--ask prices once contract design and trading costs are incorporated. The companion evidence shows that AMU is not evenly distributed across the option market. It is larger for ETH than BTC, larger away from the money than at the money, larger in short maturities than in longer maturities, and more visible during market stress. These empirical patterns are exactly the patterns one would expect if AMU were partly measuring a liquidity premium rather than a pure pricing mistake. The present paper therefore asks a different question: can the level and time evolution of AMU be interpreted as a liquidity premium that reflects the cost of enforcing the law of one price in segmented crypto derivative markets?

The motivating empirical fact is a sharp decline in PCP-breaking opportunities around the institutionalization of crypto markets in 2024. The year is an attractive candidate for a market-structure break. The U.S. Securities and Exchange Commission approved the listing and trading of spot Bitcoin exchange-traded products in January 2024, which created a regulated, brokerage-native wrapper for Bitcoin exposure \citep{sec_spot_bitcoin_etp_2024}. Spot Ether exchange-traded products followed in 2024, broadening the regulated-access channel beyond Bitcoin \citep{sec_spot_ether_etp_2024}. In the European Union, the Markets in Crypto-Assets Regulation (MiCA) became applicable in stages, with stablecoin-related provisions from June 2024 and broader crypto-asset-service-provider rules from December 2024 \citep{esma_mica_2024}. Around the same period, accounting, custody, exchange surveillance, and institutional execution channels also improved. These events did not eliminate crypto risk, but they changed the set of investors for whom arbitrage could be legally, operationally, and institutionally deployed.

Our central hypothesis is therefore not that crypto traders suddenly learned put--call parity. Rather, the hypothesis is that the no-arbitrage band narrowed. Before institutionalization, many PCP deviations may have been visible but non-bankable: large enough to observe in data but too costly, too operationally fragile, or too institutionally ambiguous to exploit at scale. After regulatory and market-structure improvements, the same deviations became actionable to a larger and better-capitalized set of arbitrageurs. This converts the empirical question from ``does parity hold?'' to ``how wide is the arbitrage band, and when does it compress?''

The contribution of the paper is threefold. First, we reinterpret AMU as a liquidity premium associated with the enforcement of PCP in crypto options. Second, we propose a threshold-arbitrage model in which parity deviations persist inside an endogenous no-arbitrage band and mean-revert faster outside it. Third, we connect the crypto evidence to mature-market episodes in which theoretically tight arbitrage relationships fail or tighten because of funding, balance-sheet, and liquidity constraints. These include covered-interest-parity deviations after the global financial crisis \citep{du2018deviations}, margin-based deviations from the law of one price \citep{garleanu2011margin}, ETF arbitrage under liquidity mismatch \citep{pan2017etf}, and the broader limits-of-arbitrage literature \citep{shleifer1997limits}.

The remainder of the paper is organized as follows. Section 2 reviews the related literature. Section 3 defines AMU as an arbitrage-band and liquidity-premium measure. Section 4 presents a threshold model of parity enforcement. Section 5 summarizes the current AMU evidence. Section 6 describes the proposed data and empirical design. Section 7 discusses historical analogies and identification. Section 8 concludes.

\section{Related literature}

\subsection{Put--call parity and crypto derivative efficiency}

Testing put--call parity has a long tradition in equity, index, and currency option markets. Early work shows that apparent parity violations are sensitive to bid--ask spreads, execution timing, short-sale constraints, taxes, and transaction costs \citep{klemkosky1979put, nisbet1992put, kamara1995daily, el1998put, cremers2010deviations}. Crypto option markets create additional complications because contracts are often inverse or quanto-style instruments, collateral and settlement currencies may differ from payoff currencies, and the underlying reference index is not always directly tradable \citep{alexander2023crypto, lucic2024valuation}.

Recent crypto work increasingly treats derivative efficiency as a joint market-structure problem. \citet{alexander2024arbitrage} develop arbitrage tests linking Bitcoin and Ether options with perpetual futures and find that crypto derivative markets have become more efficient over time, especially for longer-maturity options. \citet{felfoldi2024put} study high-frequency put--call parity in Binance Ether options and emphasize that violations become much smaller once synchronization and trading costs are handled carefully. Our companion AMU paper extends this logic to inverse options and matching futures on OKX and Deribit, constructing a directly comparable measure of investor-facing arbitrage profits in basis points.

\subsection{Liquidity premia and the limits of arbitrage}

The liquidity-premium interpretation of AMU is motivated by the idea that arbitrage is costly and capital constrained. \citet{amihud1986asset} show that bid--ask spreads affect asset pricing because investors require compensation for illiquidity. \citet{shleifer1997limits} emphasize that textbook arbitrage is typically not riskless or capital-free in practice. \citet{garleanu2011margin} formalize how margin requirements and funding constraints generate deviations from the law of one price. These ideas are directly relevant to crypto derivatives, where collateral denomination, margin rules, exchange credit risk, and stablecoin conversion all affect the cost of enforcing parity.

The closest mature-market analogy is covered interest parity (CIP). In frictionless international money markets, CIP should bind tightly. Nevertheless, post-crisis bank balance-sheet constraints and dollar funding pressures generated persistent CIP deviations \citep{du2018deviations}. The crypto case studied here is the mirror image: regulatory clarity and institutional access may have reduced, rather than increased, the effective cost of arbitrage capital, leading to a compression of parity deviations.

ETF pricing provides a second analogy. ETF prices are kept close to net asset value by authorized participants who create and redeem shares. However, when the underlying assets are illiquid or arbitrageurs face balance-sheet constraints, ETF premiums and discounts can persist \citep{pan2017etf, todorov2021bond}. The 2024 introduction of U.S. spot Bitcoin ETFs is particularly relevant because it created a regulated arbitrage and custody channel between traditional brokerage infrastructure and digital asset exposure.

\section{AMU as a liquidity premium}

Let $X_{i,t}$ denote the quoted put--call parity wedge for option pair $i$ at time $t$, expressed in basis points of the underlying notional after contract-specific conversion into a common currency. In a standard textbook setting, $X_{i,t}=0$. In the AMU framework, however, the relevant object is the positive part of the executable arbitrage profit after bid--ask prices and fees:

\begin{equation}
\mathrm{AMU}_{i,t}=\left[\frac{\Pi_{i,t}}{N_{i,t}}\right]^+\times 10{,}000,
\end{equation}

where $\Pi_{i,t}$ is the best available PCP-arbitrage profit across the implementable arbitrage paths and $N_{i,t}$ is the option contract's underlying notional. In the companion paper, $\Pi_{i,t}$ is computed from synchronized option, futures, spot, and stablecoin quotes using the correct inverse-contract payoff structure.

The present paper interprets AMU as a lower bound on the liquidity premium associated with the enforcement of parity. It captures the part of the quoted bid--ask structure that leaves a risk-free profit to a market maker who passively joins the top of the order book and completes the arbitrage cycle once filled. Thus AMU is not merely a pricing error. It is a measure of how much compensation the marginal liquidity supplier can extract because competitors do not immediately close the spread.

The no-arbitrage restriction should therefore be written as a band:

\begin{equation}
|X_{i,t}| \leq B_{i,t},
\end{equation}

where the band $B_{i,t}$ includes observable and latent implementation costs:

\begin{equation}
B_{i,t}=\tau_{i,t}+s_{i,t}+m_{i,t}+f_{i,t}+\ell_{i,t}+\kappa_{i,t}.
\end{equation}

Here $\tau_{i,t}$ denotes explicit fees, $s_{i,t}$ bid--ask and market-impact costs, $m_{i,t}$ margin and collateral costs, $f_{i,t}$ funding and stablecoin-conversion costs, $\ell_{i,t}$ legal and regulatory uncertainty, and $\kappa_{i,t}$ counterparty, custody, and operational costs. The key empirical claim is that the post-2024 market experienced a decline in several components of $B_{i,t}$.

A reduced-form representation is:

\begin{equation}
B_{i,t}=\alpha_0+\alpha_1\mathrm{Spread}_{i,t}+\alpha_2\mathrm{Depth}^{-1}_{i,t}+\alpha_3\mathrm{Funding}_{t}+\alpha_4\mathrm{Stress}_{t}+\alpha_5\mathrm{Regime}_{t}+u_{i,t},
\end{equation}

where $\mathrm{Regime}_{t}$ is a post-institutionalization indicator or a sequence of event dummies. A negative $\alpha_5$ is the empirical signature of liquidity-premium compression.

\section{A threshold model of parity enforcement}

The simplest dynamic model treats parity deviations as persistent inside the arbitrage band and rapidly mean-reverting outside it:

\begin{equation}
X_{i,t+1}=\begin{cases}
\rho_{0}X_{i,t}+\varepsilon_{i,t+1}, & |X_{i,t}|\leq B_{i,t},\\
\rho_{1}X_{i,t}+\varepsilon_{i,t+1}, & |X_{i,t}|>B_{i,t},
\end{cases}
\qquad |\rho_1|<|\rho_0|.
\end{equation}

Inside the band, arbitrageurs do not trade aggressively because the wedge is too small after implementation costs. Outside the band, arbitrage capital enters and compresses the deviation. A market-structure transition can then be represented as a break in the threshold:

\begin{equation}
B_{i,t}=B_{i}^{\mathrm{pre}}\mathbf{1}_{t<\tau}+B_{i}^{\mathrm{post}}\mathbf{1}_{t\geq\tau}, \qquad B_{i}^{\mathrm{post}}<B_{i}^{\mathrm{pre}}.
\end{equation}

This model has two testable implications. First, the distribution of $|X_{i,t}|$ and AMU should shift left after the institutionalization events. Second, conditional on large deviations, mean reversion should become faster because more capital can be deployed to enforce parity.

A more structural but still tractable version lets the band depend on arbitrage capital:

\begin{equation}
B_{i,t}=\frac{\phi_{i,t}}{A_t},
\end{equation}

where $\phi_{i,t}$ is the per-unit implementation cost and $A_t$ is effective arbitrage capacity. Regulatory clarity and ETF-linked market integration lower $\phi_{i,t}$ and increase $A_t$, both of which narrow the no-arbitrage band. This formulation captures the central mechanism in one sentence: the collapse of PCP violations reflects a transition from visible but non-bankable arbitrage to institutionally enforceable arbitrage.

\section{Evidence from the current AMU results}

The current AMU results provide the initial empirical foundation for the liquidity-premium interpretation. Although the available text in the companion paper currently focuses on 2025--early 2026 rather than the full 2020--2026 transition, the cross-sectional patterns already point toward liquidity as the organizing variable. BTCUSD options show lower AMU than ETHUSD options; at-the-money options show lower AMU than away-from-the-money options; longer-dated contracts appear more efficient than the shortest maturities; and AMU rises during periods of market stress. These patterns are difficult to reconcile with a story in which PCP violations are random pricing errors. They are more naturally interpreted as compensation left to liquidity suppliers in states where arbitrage is harder to implement.

The purpose of the longer historical extension is therefore not simply to add more observations. It is to turn these cross-sectional liquidity facts into a time-series test of market integration. If the no-arbitrage band narrowed around 2024, the same patterns should be visible in both dimensions: illiquid option segments should have high AMU at any point in time, and the market as a whole should display lower AMU after the arrival of deeper, more regulated, and more institutionally usable arbitrage channels.

\section{Data and empirical design}

\subsection{Core data}

The core sample should use Deribit and OKX option and futures order-book data from Tardis, ideally from 2020 through February 2026. The companion AMU paper currently analyzes approximately one year of OKX and Deribit data from 2025 to early 2026. The present paper should extend the same AMU calculation backward to capture the pre-2024 regime and the transition around 2024.

For each exchange, underlying, strike, maturity, and timestamp, the dataset should contain synchronized best bid and ask quotes for calls, puts, and matching expiry futures. For inverse products, all legs must be converted into a common currency using the platform-specific payoff and premium conventions. Observations should be filtered to matched put--call pairs with the same strike and maturity, positive bid quotes, and a futures contract or forward proxy with matching expiry.

\subsection{Liquidity and friction variables}

The main dependent variables are $\mathrm{AMU}_{i,t}$, an indicator $1\{\mathrm{AMU}_{i,t}>0\}$, the raw parity wedge $X_{i,t}$, and the absolute wedge $|X_{i,t}|$. The main liquidity controls should include option bid--ask spreads, futures bid--ask spreads, top-of-book depth, order-book imbalance, option open interest, realized volume, relative strike $K/F$, time to expiration, implied-volatility level, and realized volatility.

Additional market-friction variables should be collected from crypto and traditional-market sources. These include perpetual-futures funding rates, stablecoin basis measures such as USDT/USD and USDC/USD deviations, exchange-level outage or stress indicators, BTC and ETH spot returns, realized volatility, ETF flows, ETF premium/discounts, and CME futures basis. ETF variables are especially useful after January 2024 because they proxy for regulated access and institutional arbitrage capacity.

\subsection{Event variables}

The baseline event dates are January 10, 2024 for U.S. spot Bitcoin ETP approval, May 2024 and July 2024 for spot Ether ETF approval and trading, June 30, 2024 for MiCA stablecoin provisions, and December 30, 2024 for the broader MiCA crypto-asset-service-provider framework. The paper should not claim that any single date caused the transition mechanically. Instead, these dates should be treated as markers in a broader institutionalization process.

\subsection{Empirical specifications}

A first-pass panel specification is:

\begin{equation}
\mathrm{AMU}_{i,t}=\beta_0+\beta_1\mathrm{Post2024}_t+\beta_2\mathrm{Liquidity}_{i,t}+\beta_3\mathrm{Stress}_t+\gamma_i+\delta_t+\varepsilon_{i,t},
\end{equation}

where $\gamma_i$ are contract-pair fixed effects and $\delta_t$ are calendar-time fixed effects at the day or hour level. A negative $\beta_1$ after controlling for spreads, depth, and volatility would suggest that the post-2024 compression is not merely a mechanical effect of local liquidity.

A second specification estimates the probability of a tradable arbitrage opportunity:

\begin{equation}
\Pr(\mathrm{AMU}_{i,t}>0)=\Lambda\left(\beta_0+\beta_1\mathrm{Post2024}_t+\beta_2\mathrm{Liquidity}_{i,t}+\beta_3\mathrm{Stress}_t+\gamma_i+\delta_t\right).
\end{equation}

A third specification estimates the threshold model directly. For a grid of candidate thresholds $B$, estimate separate persistence coefficients inside and outside the band and select the threshold minimizing the residual sum of squares. The pre- and post-2024 thresholds can then be compared by exchange, underlying, maturity bucket, and moneyness bucket.

Finally, local projections can test whether large deviations decay faster after 2024:

\begin{equation}
|X_{i,t+h}|-|X_{i,t}|=\theta_h |X_{i,t}|+\lambda_h \left(|X_{i,t}|\times \mathrm{Post2024}_t\right)+\Gamma_h Z_{i,t}+\varepsilon_{i,t+h}.
\end{equation}

A negative $\lambda_h$ indicates faster post-2024 convergence.

\section{Historical analogies}

The crypto liquidity-premium story is best understood as part of a broader class of arbitrage relationships that are exact in theory but costly in practice. Covered interest parity is the canonical example. After the global financial crisis, CIP deviations became persistent because bank balance sheets and dollar funding constraints limited arbitrage \citep{du2018deviations}. In crypto options, the sign of the historical transition may be reversed: instead of new constraints widening the band, institutionalization and regulatory clarity narrowed it.

ETF arbitrage provides a second analogy. ETF premiums and discounts are controlled by authorized participants, but the arbitrage mechanism weakens when underlying assets are illiquid or AP balance sheets are constrained \citep{pan2017etf, todorov2021bond}. The Bitcoin ETF approvals matter in this context not only because they attracted flows, but because they inserted Bitcoin exposure into a familiar creation-redemption, custody, and brokerage architecture.

The on-the-run/off-the-run Treasury spread is a third analogy. The spread between newly issued and older Treasury securities has often been interpreted as a liquidity premium. Convergence trades can narrow the spread, but only when arbitrageurs have sufficient funding and balance sheet capacity. Similarly, the CDS--bond basis measures a deviation between synthetic and cash credit exposure and is known to vary with funding, margin, and capital costs \citep{garleanu2011margin, duffie2010presidential}.

These analogies suggest that the relevant object is not the existence of a no-arbitrage formula, but the institutional mechanism that enforces it. The AMU measure is useful precisely because it measures the quoted compensation left to the marginal liquidity supplier when that mechanism is incomplete.

\section{Expected findings and interpretation}

The extended 2020--2026 version of the paper should aim to establish four empirical facts. First, AMU should be materially higher in the pre-2024 period than in the post-2024 period, especially for the same segments that already look least liquid in the current AMU results: ETH, short-dated options, out-of-the-money options, and lower-depth quote states. Second, the estimated arbitrage threshold should fall after 2024. Third, large parity deviations should mean-revert more quickly after 2024. Fourth, the 2024 compression should remain visible after controlling for local liquidity variables, suggesting that institutionalization affected the market beyond ordinary bid--ask tightening.

The interpretation should remain cautious. The 2024 transition coincides with many simultaneous changes: higher trading volume, improved market-making technology, exchange competition, ETF approvals, accounting changes, custody improvements, and regulatory clarification. The paper should therefore avoid a monocausal claim. The stronger and more defensible statement is that AMU behaves like a liquidity-premium measure whose compression coincides with, and is plausibly supported by, the institutionalization of crypto markets.

\section{Conclusion}

This paper proposes a liquidity-premium interpretation of put--call parity violations in cryptocurrency option markets. Instead of treating PCP deviations as simple failures of rational pricing, we interpret them as the observable residue of an arbitrage band generated by liquidity, funding, collateral, custody, operational, and regulatory frictions. The AMU measure provides a natural empirical proxy for this band because it captures the positive risk-free arbitrage profit left in quoted bid--ask prices after contract-specific conversion and trading costs.

The proposed threshold model formalizes the idea that parity deviations can persist inside a no-arbitrage band and are eliminated only when they exceed the marginal cost of arbitrage. The institutionalization of crypto markets around 2024 plausibly narrowed this band by reducing legal uncertainty, improving access to regulated exposure, deepening arbitrage capital, and strengthening the market infrastructure needed to enforce cross-market pricing relationships.

By linking AMU to the literature on liquidity premia and limits to arbitrage, the paper reframes crypto option inefficiency as a familiar financial-market phenomenon. Crypto markets are unusual in contract design and settlement mechanics, but the deeper economic mechanism is classical: the law of one price is enforced only when arbitrage capital is willing and able to act.

\bibliography{bibl}

\end{document}